\subsection{Limitations and threats to validity}
\label{sec:limitations}

This section is organised around the nine pitfalls of \emph{Chasing Shadows}
\cite{evertz2026chasing}, with two checks added that the survey does not have and that both caught
errors here. It begins with damage rather than with hazards. One study is withdrawn, the
\fact{drift-n}-sample temporal-drift analysis of Section~\ref{sec:detector}, which was completed and
had entered our claim ledger as a novel result; it is not restated anywhere in this paper.
Table~\ref{tab:pitfalls} gives our exposure to each pitfall with the basis for the verdict, and four
of the nine need more than a row.

\begin{table}[!htbp]
\caption{Our exposure to each of the nine pitfalls, with the basis for the verdict.}\label{tab:pitfalls}
\centering\small
\begin{tabular}{@{}
  >{\raggedright\arraybackslash}p{0.20\linewidth}
  >{\raggedright\arraybackslash}p{0.12\linewidth}
  >{\raggedright\arraybackslash}p{0.60\linewidth}@{}}
\toprule
pitfall & verdict & basis \\
\midrule
P1 data poisoning & \texttt{CLEAR} & nothing is trained; retrieval corpora are curated and their provenance documented \\
P2 label inaccuracy & \texttt{CLEAR} & ground truth is MITRE-curated and no LLM scored any result; no human analyst rated any report either \\
P3 data leakage & \texttt{PARTIAL} & one instance found by us, disclosed and mitigated, but never systematically audited; the model's training data is unknown to us and memorisation was not probed \\
P4 model collapse & \texttt{CLEAR} & augmenting the long-term memory corpus with LLM-generated cases was considered and refused \\
P5 spurious correlations & \texttt{PARTIAL} & two were found and are reported; no systematic perturbation testing was performed \\
P6 context truncation & \texttt{EXPOSED} & truncation is designed in throughout and the counters now exist, but the distribution over a full cohort does not \\
P7 prompt sensitivity & \texttt{PARTIAL} & a structured prompt-variation study exists, but production prompts are fixed and were not varied per model \\
P8 surrogate fallacy & \texttt{PARTIAL} & every local result comes from one model on one machine, and a second endpoint has now run to completion against it \\
P9 model ambiguity & \texttt{PARTIAL} & model and engine are pinned by digest, revision and commit, but the commit was reconstructed rather than recorded \\
\bottomrule
\end{tabular}
\end{table}

On P6, truncation enters through chunking at function boundaries, a per-call evidence cap, a
\fact{react-step-budget}-step ReAct budget, an \fact{judge-output-cap}-token verdict cap and a
\setting{400}-char hint cap, and two findings belong here rather than in the row. The enumeration of
truncation sites was itself incomplete: a \fact{graph-edge-cap}-edge cap on the call-graph fetch was
discovered only while instrumenting a different measurement, and it silently truncates
\fact{graph-edge-cap-hit} of \fact{hint-total} samples, so we describe how our list was built rather
than presenting it as exhaustive. And bounds compose, as M4 shows; removing a hint made the judge
overrun a \setting{600}\,s ceiling and return an empty bundle
\fact{empty-bundles-local}/\fact{empty-bundles-pairs} times.

P8 is the pitfall this work is most exposed to and the one where the evidence has moved. A
\fact{frontier-params-b}B reasoning model has now run to completion on the same fixtures, repeats,
prompt and token budget and did not separate from the local model
(Section~\ref{sec:scale}); a \fact{series-param-span}$\times$ parameter advantage producing no
measurable difference is evidence against the pitfall's usual worry rather than merely an absence of
evidence for it. That null is confounded by the reasoning flag and the confound cannot be removed on
that endpoint, and what is being confounded is not small, the same flag being worth
\fact{reasoning-flag-worth-low} and \fact{reasoning-flag-worth-high} F1 on two separate models.
Quantisation is no longer unmeasured under P8 either, the vendor's full-precision hosting of the
same weights scoring \fact{quantisation-delta} F1 against our 3-bit deployment at
n=\fact{quantisation-pairs}, an interval containing zero and a point estimate pointing away from the
direction the threat assumes; two serving stacks differ alongside the precision, so this bounds the
deployment rather than quantisation alone. What P8 cannot close is the size series, for a measured
rather than a budgetary reason: testing a parameter-size trend needs three arms at three sizes
matched on the flag that outweighs size, the one endpoint honouring the flag hosts the same 35B
model we already run, and the one above 35B does not honour it. Coverage also remains, the second
model having been measured on five synthetic fixtures rather than on the malware cohort, which the
free tier's \setting{50} requests/day makes a two-day run or a paid one. P9 is \texttt{PARTIAL} for
one reason: the running binary reports its version as \texttt{unknown} and the engine commit was
recovered from a vendored copy of the sources rather than read off the artefact
(Section~\ref{models-endpoints-and-serving-configuration}). Build provenance must be captured at
build time; ours was reconstructed, and we say so.

The two checks the survey does not include are described as practices in Section~\ref{sec:method}
and both were added after they caught errors. All nine pitfalls concern the relationship between
claims and experiments, and what diffing the claim register against the evidence log caught sat one
level below that, needing no literature search but only two of our own documents read against each
other: a project disciplined enough to keep both a findings log and a claim ledger has, by that
fact, created the conditions for the two to diverge. The second is the output-cardinality check of
Section~\ref{sec:detector}, which found every boundary defect and which we now report alongside
results at \fact{cardinality-distinct-graphs} distinct call-graph sizes across
\fact{cardinality-samples} samples.

A scored component that derives its output from the same source as the labels is scoring a
tautology, so this was tested component by component rather than assumed. Ground truth is the MITRE
\texttt{malware uses attack-pattern} relationship set, built from the public ATT\&CK STIX
distribution, so it comes from MITRE and from nothing this project produces. Three components are
clear. The Sigma layer reads the technique identifier from each rule's own \texttt{attack.tNNNN}
tag, and those rules are community-authored detections with hand-written tags, so Sigma shares
ATT\&CK's vocabulary with the labels and not their source. The sandbox baseline, which is the one
that would have mattered most, takes its identifiers from class attributes hand-written in each
signature module and mapped through the sandbox's own table, consulting no family attribution
anywhere, which is precisely what makes it a valid baseline rather than a second view of the answer.
And the retrieval corpora do not overlap the evaluation cohorts, checked by digest rather than by
assertion: of the \fact{case-corpus-n} cases in the ATT\&CK case corpus, none shares a sha256 with
the \fact{cohort-n}-sample dynamic cohort or with the \fact{drift-n}-sample drift manifest. One
construct problem survives and it is ours: the case corpus is labelled with the technique
identifiers the pipeline itself previously attributed, so scoring the index against those labels
measures agreement with its own upstream, which is why that number is relabelled as self-consistency
in Section~\ref{sec:mechanisms} and the conclusion rests on the production figure beside it. There
is also one shared vocabulary, named rather than hidden: the catalogue defining valid technique
identifiers is the same file that bounds the label universe, so the invalid-identifier rate is
measured against the catalogue the labels are drawn from.

Any claim about what the sandbox contributes has to survive one measurement of what the sandbox
reports. Across the \fact{evidence-reports} archived analyses there are
\fact{network-domains-distinct} distinct network domains and \fact{network-domains-ubiquitous} of
them appear in every single sample, the analysis VM's own vendor telemetry, present in each capture
whatever was detonated. Per sample, between \fact{network-ubiquitous-min} and
\fact{network-ubiquitous-max} of observed domains are cohort-ubiquitous, median
\fact{network-ubiquitous-median}, so two-thirds of the network evidence in a dynamic arm is the
instrument describing itself and an effect attributed to malware behaviour may be partly an effect
of the analysis VM's idle traffic. The same measurement bounds two Layer-0 sources from the other
direction: the LOLBin and DGA layers produced a claim on none of the \fact{layer0-declined-of}
reports archived when that ablation ran, while being fed a median of
\fact{layer0-declined-median-calls} recorded API calls and \fact{layer0-declined-domains} domains
respectively, so neither appears as an arm in the cascade ablation. An earlier version of this
measurement, on the \fact{ubiquity-interim-reports} analyses that survived before the cohort was
recovered, gave \fact{ubiquity-interim-domains} domains with \fact{ubiquity-interim-ubiquitous}
ubiquitous at a median share of \fact{ubiquity-interim-median}; more than doubling the cohort moved
the figure by three points and left the conclusion where it was.

The machine a measurement runs on is part of the instrument, which is a threat we did not anticipate
and met late. The working set does not fit a \fact{host-ram-gib} GiB host alongside an interactive
session, with the model server holding \textasciitilde16 GB, the analysis worker
\textasciitilde8 GB and the disassembly container up to \fact{ghidra-mem-limit-gb} GB. During a
paired ablation the swap file was exhausted and the model server had \fact{llama-paged-out-gb} GB of
its own address space paged out; an arm then exceeded a \fact{salvage-budget-seconds}-second budget
on a prompt trimmed to \fact{static-max-chars} characters. Whether that arm failed because of the
pipeline or because of the host cannot be determined, because we recorded the pipeline's outputs and
not the host's state, so the affected arms are reported as \texttt{unattributable} rather than as
failures and the ablation as incomplete rather than as a null. Three things follow. Wall-clock
bounds are host-dependent measurements wearing the costume of a system property, so any result here
that turns on a timeout is a statement about this pipeline on this machine under whatever load it
had. The retention rule adopted after the first failure did not cover this one, because it was
written about the object under study and what went unrecorded was the environment the study ran in.
And a screen for this can only be applied forward, since arms already collected without host state
cannot be re-screened, which is the same structural problem that withdrew the drift study appearing
in a new variable. A postscript, because it took three attempts to see: two of those interruptions
we attributed to memory pressure and treated with progressively better limits, a floor, then a
strike counter, then an allowance for declared allocations, then marking the model server as the
kernel's preferred OOM victim. All four were correct and none was the cause. The harness had started
the server as a child process, so it ran inside the cgroup of the terminal that launched it and its
memory was accounted to that scope; the kernel's log named the scope plainly and we did not read it
until the third failure. Every fix chose which process the kernel would kill, and the defect was
whose accounting it was killed inside.

Three objections remain that the work does not close, and they are stated rather than answered. The
first is that the judge findings rest on four calls, because the intercept patches the reconciliation
seam and the text-fallback builder and therefore requires a live judge call per observation. We
concede the sample size and defend the design, which are different things. On the working path the
finding is deductive rather than statistical, since the reconciliation step restores the cascade's
set by construction and the \fact{judge-mechanism-arms}-arm mechanism check confirms the bundle
equals the cascade set on every arm, with the judge contributing one of its own on
\fact{judge-mechanism-added} of \fact{judge-mechanism-samples} samples, an upper bound of
\fact{judge-mechanism-upper-bound} on the per-sample rate. On the failing path it is an existence
claim: uncorroborated identifiers do reach the analyst, demonstrated on four calls against an
uncapped control whose raw text provably contains none. Neither claim is a rate, and what would
settle it is the same intercept over the \fact{baseline-n}-sample cohort, which needs no new
instrumentation, only the host. The second is that the detector's own evidence is retrospective: it
found four defects on batches we already had reason to distrust, and the prospective test, reporting
output cardinality beside every batch for a year and counting how many defects it catches before
anything else does, is the experiment and we have not run it. The third is that seven defects in
one's own code is a post-mortem rather than a result. We do not claim the category. Every mechanism
was counter-searched in an adjacent field's vocabulary before being claimed and the search demoted
more than it confirmed, so what we assert is narrower than a class: a setting where the artefact of
such a failure is a measurement that is wrong and looks right, and a price paid rather than
hypothesised.

All three were blocked by the same thing, a host on which the model server does not fit alongside
anything else. Lifting an earlier objection shows what that takes: the blocker was never the method
but that the server's resident set climbs with cumulative requests and does not come back down, at
\fact{llama-rss-fresh-gb} GB loaded and \fact{llama-rss-sustained-gb} GB after
\fact{llama-rss-sustained-minutes} minutes, unreclaimable, so a long run walks the host into its own
memory guard. Cycling the server from the checkpoint between stretches of work turned a run that
could not finish into one that did, which is a change to the method and not to the machine. That is
what closed the objection that the equal-budget arms ran only on \fact{fixture-clusters} synthesised
fixtures whose evidence is in bijection with their answer key, a corpus on which a deterministic
regular expression over the generating dictionary scores \fact{fixture-extractor-f1} on
\fact{fixture-extractor-n} of \fact{fixture-extractor-n}: two of the three arms were re-run on real
binaries and the answer changed sign (Section~\ref{sec:consensus}). The frontier comparison and the
confidence calibration still rest on that corpus and are reported as a pilot accordingly.
